\ifdefined\gkver\else\def\gkver{student}\fi
\documentclass[\gkver]{gaokaozhenti}
\begin{document}

\chapter{2013年北京理综}
%% paperType: 理综物理部分

\begin{enumerate}
\item
%% number: 13
%% paperName: 2013年北京理综第 13 题 6 分
%% typeId: 1
%% type: 单选题
%% chapter: 气体性质
%% point: 分子动理论
%% method: 无
%% score: 6
%% degree: 850
%% duplicateId: 0
%% body:
下列说法正确的是\xzanswer{A}
\fourchoices[answer=A]
{液体中悬浮的微粒的无规则运动称为布朗运动}
{液体分子的无规则运动称为布朗运动}
{物体从外界吸收热量，其内能一定增加}
{物体对外界做功，其内能一定减少}
%% answer: A
%% memo:
\memoanswer{
本题原卷及材料中未提供详解，待补充。
}

\item
%% number: 14
%% paperName: 2013年北京理综第 14 题 6 分
%% typeId: 1
%% type: 单选题
%% chapter: 光学
%% point: 光的折射
%% method: 无
%% score: 6
%% degree: 650
%% duplicateId: 0
%% body:
如图所示，一束可见光射向半圆形玻璃砖的圆心 $O$，经折射后分为两束单色光 $a$ 和 $b$。下列判断正确的是\xzanswer{B}
\onepicture[w=4.5cm, align=right]{figs/14.png}
\fourchoices[answer=B]
{玻璃对 $a$ 光的折射率小于对 $b$ 光的折射率}
{$a$ 光的频率大于 $b$ 光的频率}
{在真空中 $a$ 光的波长大于 $b$ 光的波长}
{$a$ 光光子能量小于 $b$ 光光子能量}
%% answer: B
%% memo:
\memoanswer{
本题原卷及材料中未提供详解，待补充。
}

\item
%% number: 15
%% paperName: 2013年北京理综第 15 题 6 分
%% typeId: 1
%% type: 单选题
%% chapter: 机械波
%% point: 波的图象
%% method: 无
%% score: 6
%% degree: 650
%% duplicateId: 0
%% body:
一列沿 $x$ 轴正方向传播的简谐机械横波，波速为 $4\Ums$。某时刻波形如图所示，下列说法正确的是\xzanswer{D}
\onepicture[w=6cm, align=right]{figs/15.png}
\fourchoices[answer=D]
{这列波的振幅为 $4\Ucm$}
{这列波的周期为 $1\Us$}
{此时 $x=4\Um$ 处质点沿 $y$ 轴负方向运动}
{此时 $x=4\Um$ 处质点的加速度为 0}
%% answer: D
%% memo:
\memoanswer{
本题原卷及材料中未提供详解，待补充。
}

\item
%% number: 16
%% paperName: 2013年北京理综第 16 题 6 分
%% typeId: 1
%% type: 单选题
%% chapter: 力物体的平衡
%% point: 物体系平衡
%% method: 整体与隔离
%% score: 6
%% degree: 650
%% duplicateId: 0
%% body:
倾角为 $\alpha$、质量为 $M$ 的斜面体静止在水平桌面上，质量为 $m$ 的木块静止在斜面体上。下列结论正确的是\xzanswer{D}
\onepicture[w=5.5cm, align=right]{figs/16.png}
\fourchoices[answer=D]
{木块受到的摩擦力大小是 $mg\cos\alpha$}
{木块对斜面体的压力大小是 $mg\sin\alpha$}
{桌面对斜面体的摩擦力大小是 $mg\sin\alpha\cos\alpha$}
{桌面对斜面体的支持力大小是 $(M+m)g$}
%% answer: D
%% memo:
\memoanswer{
本题原卷及材料中未提供详解，待补充。
}

\item
%% number: 17
%% paperName: 2013年北京理综第 17 题 6 分
%% typeId: 1
%% type: 单选题
%% chapter: 电磁感应
%% point: 切割磁感线电动势
%% method: 无
%% score: 6
%% degree: 650
%% duplicateId: 0
%% body:
如图，在磁感应强度为 $B$、方向垂直纸面向里的匀强磁场中，金属杆 MN 在平行金属导轨上以速度 $v$ 向右匀速滑动，MN 中产生的感应电动势为 $E_{1}$；若磁感应强度增为 $2B$，其他条件不变，MN 中产生的感应电动势变为 $E_{2}$。则通过电阻 $R$ 的电流方向及 $E_{1}$ 与 $E_{2}$ 之比 $E_{1}:E_{2}$ 分别为\xzanswer{C}
\onepicture[w=5.5cm, align=right]{figs/17.png}
\fourchoices[answer=C]
{$c\to a$，$2:1$}
{$a\to c$，$2:1$}
{$a\to c$，$1:2$}
{$c\to a$，$1:2$}
%% answer: C
%% memo:
\memoanswer{
本题原卷及材料中未提供详解，待补充。
}

\item
%% number: 18
%% paperName: 2013年北京理综第 18 题 6 分
%% typeId: 1
%% type: 单选题
%% chapter: 电场
%% point: 带电粒子的圆周运动
%% method: 无
%% score: 6
%% degree: 650
%% duplicateId: 0
%% body:
某原子电离后其核外只有一个电子，若该电子在核的静电力作用下绕核做匀速圆周运动，那么电子运动\xzanswer{C}
\fourchoices[answer=C]
{半径越大，加速度越大}
{半径越小，周期越大}
{半径越大，角速度越小}
{半径越小，线速度越小}
%% answer: C
%% memo:
\memoanswer{
本题原卷及材料中未提供详解，待补充。
}

\item
%% number: 19
%% paperName: 2013年北京理综第 19 题 6 分
%% typeId: 1
%% type: 单选题
%% chapter: 曲线运动
%% point: 平抛运动
%% method: 无
%% score: 6
%% degree: 650
%% duplicateId: 0
%% body:
在实验操作前应该对实验进行适当的分析。研究平抛运动的实验装置示意如图。小球每次都从斜槽的同一位置无初速度释放，并从斜槽末端水平飞出。改变水平板的高度，就改变了小球在板上落点的位置，从而可描绘出小球的运动轨迹。某同学设想小球先后三次做平抛，将水平板依次放在如图 1、2、3 的位置，且 1 与 2 的间距等于 2 与 3 的间距。若三次实验中，小球从抛出点到落点的水平位移依次为 $x_{1}$、$x_{2}$、$x_{3}$，机械能的变化量依次为 $\Delta E_{1}$、$\Delta E_{2}$、$\Delta E_{3}$，忽略空气阻力的影响，下面分析正确的是\xzanswer{B}
\onepicture[w=6cm, align=right]{figs/19.png}
\fourchoices[answer=B]
{$x_{2}-x_{1}=x_{3}-x_{2}$，$\Delta E_{1}=\Delta E_{2}=\Delta E_{3}$}
{$x_{2}-x_{1}>x_{3}-x_{2}$，$\Delta E_{1}=\Delta E_{2}=\Delta E_{3}$}
{$x_{2}-x_{1}>x_{3}-x_{2}$，$\Delta E_{1}<\Delta E_{2}<\Delta E_{3}$}
{$x_{2}-x_{1}<x_{3}-x_{2}$，$\Delta E_{1}<\Delta E_{2}<\Delta E_{3}$}
%% answer: B
%% memo:
\memoanswer{
本题原卷及材料中未提供详解，待补充。
}

\item
%% number: 20
%% paperName: 2013年北京理综第 20 题 6 分
%% typeId: 1
%% type: 单选题
%% chapter: 光学
%% point: 光电效应
%% method: 无
%% score: 6
%% degree: 450
%% duplicateId: 0
%% body:
以往我们认识的光电效应是单光子光电效应，即一个电子在极短时间内只能吸收到一个光子而从金属表面逸出。强激光的出现丰富了人们对于光电效应的认识，用强激光照射金属，由于其光子密度极大，一个电子在极短时间内吸收多个光子成为可能，从而形成多光子电效应，这已被实验证实。

光电效应实验装置示意如图。用频率为 $\nu$ 的普通光源照射阴极 $K$，没有发生光电效应。换用同样频率为 $\nu$ 的强激光照射阴极 $K$，则发生了光电效应；此时，若加上反向电压 $U$，即将阴极 $K$ 接电源正极，阳极 $A$ 接电源负极，在 $KA$ 之间就形成了使光电子减速的电场，逐渐增大 $U$，光电流会逐渐减小；当光电流恰好减小到零时，所加反向电压 $U$ 可能是下列的（其中 $W$ 为逸出功，$h$ 为普朗克常量，$e$ 为电子电量）\xzanswer{B}
\onepicture[w=4.5cm, align=right]{figs/20.png}
\fourchoices[answer=B]
{$U=\frac{h\nu}{e}-\frac{W}{e}$}
{$U=\frac{2h\nu}{e}-\frac{W}{e}$}
{$U=2h\nu-W$}
{$U=\frac{5h\nu}{2e}-\frac{W}{e}$}
%% answer: B
%% memo:
\memoanswer{
本题原卷及材料中未提供详解，待补充。
}

\item
%% number: 21
%% paperName: 2013年北京理综第 21 题 18 分
%% typeId: 4
%% type: 实验题
%% chapter: 电路
%% point: 电路实验
%% method: 无
%% score: 18
%% degree: 450
%% duplicateId: 0
%% body:
某同学通过实验测定一个阻值约为 $5\UO$ 的电阻 $R_{x}$ 的阻值。
\begin{enumerate}
	\item
	现有电源（$4\UV$，内阻可不计）、滑动变阻器（$0\sim50\UO$，额定电流 $2\UA$），开关和导线若干，以及下列电表：

	（A）电流表（$0\sim3\UA$，内阻约 $0.025\UO$）

	（B）电流表（$0\sim0.6\UA$，内阻约 $0.125\UO$）

	（C）电压表（$0\sim3\UV$，内阻约 $3\UkO$）

	（D）电压表（$0\sim15\UV$，内阻约 $15\UkO$）

	为减小测量误差，在实验中，电流表应选用\tkanswer[1.2cm]{B}，电压表应选用\tkanswer[1.2cm]{C}（选填器材前的字母）；实验电路应采用图 1 中的\tkanswer[1.2cm]{甲}（选填“甲”或“乙”）。

	\onepicture[w=10cm, align=right]{figs/21a.png}

	\item
	图 2 是测量 $R_{x}$ 的实验器材实物图，图中已连接了部分导线。请根据在（1）问中所选的电路图，补充完成图 2 中实物间的连线。

	\drawpicanswer{\onepicture[w=10cm, align=right]{figs/21b.png}}{\onepicture[w=10cm, align=right]{figs/21b_an.png}}

	\item
	接通开关，改变滑动变阻器滑片 $P$ 的位置，并记录对应的电流表示数 $I$、电压表示数 $U$。某次电表示数如图 3 所示，可得该电阻的测量值 $R_{x}=\frac{U}{I}=$\tkanswer[1.5cm]{$5.2\UO$}（保留两位有效数字）。

	\onepicture[w=10cm, align=right]{figs/21c.png}

	\item
	若在（1）问中选用甲电路，产生误差的主要原因是\tkanswer[1.2cm]{B}；若在（1）问中选用乙电路，产生误差的主要原因是\tkanswer[1.2cm]{D}。（选填选项前的字母）

	（A）电流表测量值小于流经 $R_{x}$ 的电流值

	（B）电流表测量值大于流经 $R_{x}$ 的电流值

	（C）电压表测量值小于 $R_{x}$ 两端的电压值

	（D）电压表测量值大于 $R_{x}$ 两端的电压值

	\item
	在不损坏电表的前提下，将滑动变阻器滑片 $P$ 从一端滑向另一端，随滑片 $P$ 移动距离 $x$ 的增加，被测电阻 $R_{x}$ 两端的电压 $U$ 也随之增加，下列反映 $U-x$ 关系的示意图中正确的是\tkanswer[1.2cm]{A}。

	\onepicture[w=11cm, align=right]{figs/21d.png}
\end{enumerate}
%% answer: 
\jdanswer{
\begin{enumerate}
	\item
	B、C、甲

	\item
	如图

	\item
	$5.2\UO$

	\item
	B、D

	\item
	A

\end{enumerate}
}
%% memo:
\memoanswer{
本题原卷及材料中未提供详解，待补充。
}

\item
%% number: 22
%% paperName: 2013年北京理综第 22 题 16 分
%% typeId: 6
%% type: 计算题
%% chapter: 磁场
%% point: 带电粒子在磁场中的运动
%% method: 无
%% score: 16
%% degree: 650
%% duplicateId: 0
%% body:
如图所示，两平行金属板间距为 $d$，电势差为 $U$，板间电场可视为匀强电场；金属板下方有一磁感应强度为 $B$ 的匀强磁场。带电量为 $+q$、质量为 $m$ 的粒子，由静止开始从正极板出发，经电场加速后射出，并进入磁场做匀速圆周运动。忽略重力的影响，求：
\begin{enumerate}
	\item
	匀强电场场强 $E$ 的大小；
	\item
	粒子从电场射出时速度 $v$ 的大小；
	\item
	粒子在磁场中做匀速圆周运动的半径 $R$。
\end{enumerate}
\onepicture[w=5cm, align=right]{figs/22.png}
%% answer: 
\jdanswer{
\begin{enumerate}
	\item
	$E=\frac{U}{d}$

	\item
	$v=\sqrt{\frac{2qU}{m}}$

	\item
	$R=\frac{1}{B}\sqrt{\frac{2mU}{q}}$

\end{enumerate}
}
%% memo:
\memoanswer{
（1）电场强度 $E=\frac{U}{d}$

（2）根据动能定理，有 $qU=\frac{1}{2}mv^{2}$

得 $v=\sqrt{\frac{2qU}{m}}$

（3）粒子在磁场中做匀速圆周运动时，洛伦兹力提供向心力，有

$qvB=m\frac{v^{2}}{R}$

得 $R=\frac{1}{B}\sqrt{\frac{2mU}{q}}$
}

\item
%% number: 23
%% paperName: 2013年北京理综第 23 题 18 分
%% typeId: 6
%% type: 计算题
%% chapter: 功和能
%% point: 动能定理
%% method: 无
%% score: 18
%% degree: 650
%% duplicateId: 0
%% body:
蹦床比赛分成预备运动和比赛动作。最初，运动员静止站在蹦床上；在预备运动阶段，他经过若干次蹦跳，逐渐增加上升高度，最终达到完成比赛动作所需的高度；此后，进入比赛动作阶段。

把蹦床简化为一个竖直放置的轻弹簧，弹力大小 $F=kx$（$x$ 为床面下沉的距离，$k$ 为常量）。质量 $m=50\Ukg$ 的运动员静止站在蹦床上，床面下沉 $x_{0}=0.10\Um$；在预备运动中，假定运动员所做的总功 $W$ 全部用于其机械能；在比赛动作中，把该运动员视作质点，其每次离开床面做竖直上抛运动的腾空时间均为 $\Delta t=2.0\Us$，设运动员每次落下使床面压缩的最大深度均为 $x_{1}$。取重力加速度 $g=10\Umsq$，忽略空气阻力的影响。
\begin{enumerate}
	\item
	求常量 $k$，并在图中画出弹力 $F$ 随 $x$ 变化的示意图；
	\item
	求在比赛动作中，运动员离开床面后上升的最大高度 $h_{m}$；
	\item
	借助 $F-x$ 图像可以确定弹性做功的规律，在此基础上，求 $x_{1}$ 和 $W$ 的值。
\end{enumerate}
\onepicture[w=3.5cm, align=right]{figs/23.png}
%% answer: 
\jdanswer{
\begin{enumerate}
	\item
	$k=\frac{mg}{x_{0}}=5.0\times10^{3}\UNm$，$F-x$ 图线如图

	\item
	$h_{m}=5\Um$

	\item
	$x_{1}=1.1\Um$，$W=2525\UJ=2.5\times10^{3}\UJ$

\end{enumerate}
}
%% memo:
\memoanswer{
（1）床面下沉 $x_{0}=0.10\Um$，运动员受力平衡，有 $mg=kx_{0}$，得 $k=\frac{mg}{x_{0}}=5.0\times10^{3}\UNm$，$F-x$ 图线如图。

\onepicture[w=3.5cm, align=right]{figs/23_an.png}

（2）运动员从 $x=0$ 处离开床面，开始腾空，其上升、下落时间相等，有 $h_{m}=\frac{1}{2}g\left(\frac{\Delta t}{2}\right)^{2}=5.0\Um$。

（3）参考由速度-时间图像求位移的方法，$F-x$ 图线下的面积等于弹力做的功，从 $x$ 处到 $x=0$，弹力做功 $W_{T}=\frac{1}{2}\cdot x\cdot kx=\frac{1}{2}kx^{2}$。运动员从 $x_{1}$ 处上升到最大高度 $h_{m}$ 的过程，根据动能定理，有 $\frac{1}{2}kx_{1}^{2}-mg(x_{0}+h_{m})=0$，得 $x_{1}=x_{0}+\sqrt{x_{0}^{2}+2x_{0}h_{m}}=1.1\Um$。对整个预备运动，由题设条件以及功和能的关系，有 $W+\frac{1}{2}kx_{0}^{2}=mg(h_{m}+x_{0})$，得 $W=2525\UJ=2.5\times10^{3}\UJ$。
}

\item
%% number: 24
%% paperName: 2013年北京理综第 24 题 20 分
%% typeId: 6
%% type: 计算题
%% chapter: 磁场
%% point: 洛伦兹力
%% method: 无
%% score: 20
%% degree: 450
%% duplicateId: 0
%% body:
对于同一物理问题，常常可以从宏观与微观两个不同角度进行研究，找出其内在联系，从而更加深刻地理解其物理本质。
\begin{enumerate}
	\item
	一段横截面积为 $S$、长为 $l$ 的直导线，单位体积内有 $n$ 个自由电子，电子电量为 $e$。该导线通有电流时，假设自由电子定向移动的速率均为 $v$。

	（a）求导线中的电流 $I$；

	（b）将该导线放在匀强磁场中，电流方向垂直于磁感应强度 $B$，导线所受安培力大小为 $F_{\text{安}}$，导线内自由电子所受洛伦兹力大小的总和为 $F$，推导 $F_{\text{安}}=F$。

	\item
	正方体密闭容器中有大量运动粒子，每个粒子质量为 $m$，单位体积内粒子数量 $n$ 为恒量。为简化问题，我们假定：粒子大小可以忽略；其速率均为 $v$，且与器壁各面碰撞的机会均等；与器壁碰撞前后瞬间，粒子速度方向都与器壁垂直，且速率不变。利用所学力学知识，导出器壁单位面积所受粒子压力 $f$ 与 $m$、$n$ 和 $v$ 的关系。
\end{enumerate}

（注意：解题过程中需要用到、但题目没有给出的物理量，要在解题时做必要的说明）
%% answer: 
\jdanswer{
\begin{enumerate}
	\item
	（a）设 $\Delta t$ 时间内通过导体横截面的电量为 $\Delta q$，由电流定义，有 $I=\frac{\Delta q}{\Delta t}=\frac{neSv\Delta t}{\Delta t}=neS$。

	（b）每个自由电子所受的洛伦兹力 $F_{\text{洛}}=evB$。设导体中共有 $N$ 个自由电子，$N=nSl$，导体中自由电子所受洛伦兹力大小的总和 $F=NF_{\text{洛}}=nSl\cdot evB$；由安培力公式，有 $F_{\text{安}}=IlB=neSv\cdot lB$，得 $F_{\text{安}}=F$。

	\item
	一个粒子每与器壁碰撞一次给器壁的冲量为 $\Delta I=2mv$。以器壁上的面积 $S$ 为底，以 $v\Delta t$ 为高构成柱体，由题设可知，其内的粒子在 $\Delta t$ 时间内有 $\frac{1}{6}$ 与器壁 $S$ 发生碰撞，碰壁粒子总数为 $N=\frac{1}{6}nSv\Delta t$，$\Delta t$ 时间内粒子的冲量为 $I=N\cdot 2mv=\frac{1}{3}nSmv^{2}\Delta t$。面积为 $S$ 的器壁受到粒子压力为 $F=\frac{I}{\Delta t}$，器壁单位面积所受粒子压力为 $f=\frac{F}{S}=\frac{1}{3}nmv^{2}$。

\end{enumerate}
}
%% memo:
\memoanswer{
对于同一物理问题，常常可以从宏观与微观两个不同角度进行研究，找出其内在联系，从而更加深刻地理解其物理本质。

（1）（a）设 $\Delta t$ 时间内通过导体横截面的电量为 $\Delta q$，由电流定义，有

$I=\frac{\Delta q}{\Delta t}=\frac{neSv\Delta t}{\Delta t}=neS$

（b）每个自由电子所受的洛伦兹力 $F_{\text{洛}}=evB$。设导体中共有 $N$ 个自由电子，$N=n\cdot Sl$，导体中自由电子所受洛伦兹力大小的总和

$F=NF_{\text{洛}}=nSl\cdot evB$

由安培力公式，有 $F_{\text{安}}=IlB=neSv\cdot lB$，得 $F_{\text{安}}=F$。

（2）一个粒子每与器壁碰撞一次给器壁的冲量为 $\Delta I=2mv$。如图，以器壁上的面积 $S$ 为底，以 $v\Delta t$ 为高构成柱体，由题设可知，其内的粒子在 $\Delta t$ 时间内有 $\frac{1}{6}$ 与器壁 $S$ 发生碰撞，碰壁粒子总数为

$N=\frac{1}{6}nSv\Delta t$

$\Delta t$ 时间内粒子的冲量为

$I=N\cdot 2mv=\frac{1}{3}nSmv^{2}\Delta t$

面积为 $S$ 的器壁受到粒子压力为 $F=\frac{I}{\Delta t}$，器壁单位面积所受粒子压力为

$f=\frac{F}{S}=\frac{1}{3}nmv^{2}$

\onepicture[w=3.5cm, align=right]{figs/24_an.png}

【解二】考虑单位面积，$t$ 时间内能达到容器壁的粒子所占据的体积为 $V=Svt=1\times vt$，其中粒子有均等的概率与容器各面相碰，即可能达到目标区域的粒子数为 $\frac{1}{6}nV=\frac{1}{6}nvt$，由动量定理可得 $f=\frac{\Delta P}{t}=\frac{1}{6}nvt(2\times mv)=\frac{1}{3}nmv^{2}$。
}

\end{enumerate}

\end{document}
